\chapter{CIRCONDATO DAI SOGNI}\label{circondato-dai-sogni}

\hfill\break

Le battaglie notturne tra pirati della strada e Polizia Stradale della
California nel migliore dei casi rendevano difficile viaggiare. Gli
sfarfallii peggioravano le cose. Durante gli episodi di sfarfallio le
autorità sconsigliavano qualunque tipo di viaggio non strettamente
necessario, ma questo non bastava a impedire alla gente di cercare di
raggiungere familiari e amici o, in alcuni casi, salire in macchina e
guidare fino a esaurire benzina o tempo. Avevo riempito in tutta fretta
un paio di valigie stipandoci tutto ciò che non volevo dimenticare,
compreso l'archivio di documenti che mi aveva dato Jase.

Quella notte, la tangenziale di Alvarado era intasata dal traffico e la
situazione sull'Interstatale 8 non era migliore. Avevo un sacco di tempo
per riflettere sull'assurdità di ciò che stavo provando a fare.

Precipitarmi in soccorso della moglie di un altro, della donna a cui una
volta tenevo molto più di quanto mi facesse bene. Quando chiudevo gli
occhi provando a visualizzare Diane Lawton non c'era più nessuna
rappresentazione coerente, solo un confuso montaggio di momenti e gesti.
Diane che si scostava i capelli con la mano e si appoggiava al manto di
Sant'Agostino, il suo cane. Diane che passava al fratello di nascosto un
palmare wi-fi nel capanno degli attrezzi sul cui pavimento giaceva un
tagliaerba smontato. Diane che leggeva poesie vittoriane sotto la
chiazza d'ombra di un salice e sorrideva per delle frasi che io non
avevo capito: ``\emph{L'estate matura a tutte le ore}'' o
``\emph{All'infante non è occorso}\ldots''

Diane, i cui sguardi e gesti più discreti avevano sempre sottinteso
l'amore per me, sia pure esitante, ma che era sempre stata trattenuta da
forze che non riuscivo a capire: suo padre, Jason, lo Spin. Lo Spin, che
ci aveva legato e separato, tenendoci rinchiusi in stanze adiacenti ma
senza porte, pensai.

Avevo già superato El Centro quando la radio segnalò una ``rilevante''
attività di polizia a ovest di Yuma e traffico intasato per almeno
cinque chilometri al confine di stato. Per evitare lunghe code decisi di
imboccare un raccordo secondario che attraversava il deserto vuoto a
nord e che dalla mappa sembrava promettente. Intendevo immettermi
nell'Interstatale 10 nel punto in cui questa incrociava il confine di
stato vicino a Blythe.

La strada era meno affollata ma comunque trafficata. Lo sfarfallio
faceva sembrare il mondo notturno alla rovescia, più luminoso sopra che
sotto. Ogni tanto una vena di luce particolarmente densa si contorceva
dall'orizzonte settentrionale a quello meridionale, come se si fosse
incrinata la membrana e potessimo intravedere frammenti in fiamme
dell'universo accelerato oltre lo Spin.

Pensai al telefono che avevo in tasca, il telefono di Diane, il numero
chiamato da Simon. Non potevo richiamare: non avevo il numero di Diane e
quello del ranch - sempre che si trovassero ancora là - non era
sull'elenco. Volevo soltanto che squillasse di nuovo. Lo volevo e lo
temevo.

Il traffico peggiorò di nuovo nei pressi di Palo Verde, dove la strada
si avvicinava alla statale. Ormai era passata la mezzanotte e nei tratti
liberi riuscivo al massimo a percorrere una cinquantina di chilometri
all'ora. Pensai di mettermi a dormire. Avevo bisogno di riposare. Quindi
mi decisi: mi sarei preso una pausa per il resto della notte concedendo
così al traffico il tempo di smaltirsi.

Ma non volevo dormire in macchina. Le uniche auto in sosta che avevo
visto erano state abbandonate e saccheggiate, i portabagagli spalancati
come bocche su facce sorprese.

A sud di una cittadina chiamata Ripley notai, appena visibili alla luce
dei fari, un cartello scolorito dal Sole e corroso dalla sabbia con la
scritta {ALBERGO} e una strada a doppia corsia a malapena asfaltata che
sbucava dall'autostrada. Svoltai. Cinque minuti dopo mi trovai di fronte
al cancello del complesso di edifici di quello che era - o forse era
stato un tempo - un motel, una striscia a ferro di cavallo di stanze su
due piani attorno a una piscina che sotto lo sfarfallio del cielo
sembrava vuota.

Scesi dall'auto e suonai il campanello.

Il cancello era automatico, del genere che puoi gestire a distanza di
sicurezza da un pannello di controllo, e dotato di una videocamera
grande un palmo montata su di un alto palo. La videocamera roteava per
ispezionarmi mentre un altoparlante all'altezza del finestrino crepitava
dando segni di vita. Da qualche parte, da un gabbiotto o dall'atrio del
motel, provenivano delle note musicali. Non era musica trasmessa da un
programma alla radio, sembrava riprodotta da qualche apparecchio. Poi
una voce. Brusca, metallica e ostile. «Non accettiamo ospiti stanotte.»

Pochi istanti dopo stesi il braccio e pigiai di nuovo il pulsante.

La voce ricomparve. «Che parte non ti è chiara di quello che ho detto?»

Risposi: «Posso pagare in contanti se questo cambia le cose. Non starò
lì a trattare sul prezzo.»

«Niente da fare. Mi spiace, socio.»

«Okay, aspetta solo un attimo\ldots{} senti, posso dormire in macchina
ma sarebbe un problema se venissi all'interno per essere un po' più
protetto? Parcheggiando magari sul retro, dove dalla strada non possono
vedermi?»

Pausa più lunga. Sentivo il suono di una tromba che inseguiva un
rullante. Il pezzo mi era fastidiosamente familiare.

«Scusa. Stanotte no. Per favore, vattene.»

Ancora silenzio. Altri minuti trascorsi. Da una piccola oasi di palme e
ghiaia davanti al motel, sentivo provenire il canto di un grillo. Suonai
di nuovo il campanello.

Il proprietario si rifece vivo all'istante. «Te lo devo dire, siamo
armati e alquanto incazzati. Sarebbe meglio se ti rimettessi in
viaggio.»

«\emph{Harlem Air Shaft}» dissi.

«Come?»

«Il pezzo che stai ascoltando. Ellington, vero? \emph{Harlem Air Shaft}.
Sembra proprio la sua band degli anni Cinquanta.»

Un'altra lunga pausa, sebbene l'altoparlante fosse ancora acceso. Ero
quasi certo di avere ragione nonostante non sentissi quel pezzo di Duke
Ellington da anni.

Poi la musica si interruppe, il suo filo sottilissimo tranciato a metà
battuta. «C'è qualcun altro in macchina con te?»

Tirai giù il finestrino e accesi la luce interna. La videocamera fece
una panoramica, poi puntò ancora su di me.

«Va bene» concesse. «D'accordo. Dimmi chi suona la tromba in quel pezzo
e io aziono il cancello.»

Tromba? Pensando al Duke Ellington della metà degli anni Cinquanta mi
venne in mente Paul Gonsalves, ma Gonsalves suonava il sax. C'erano
stati diversi trombettisti. Cat Anderson? Willie Cook? Era passato
troppo tempo.

«Ray Nance» abbozzai.

«No. Clark Terry. Ma suppongo che tu possa entrare lo stesso.»

ooo

Il proprietario uscì per venirmi incontro appena mi fermai davanti
all'ingresso. Era alto, sulla quarantina, indossava jeans e un'ampia
maglia a quadri. Mi osservò attentamente.

«Senza offesa,» disse, «ma la prima volta che è successo tutto
questo\ldots» indicando il cielo, lo sfarfallio che rendeva gialla la
sua pelle e di un ocra infetto le pareti in stucco. «Insomma, quando
hanno chiuso le frontiere a Blythe, è arrivata gente che si picchiava
per avere una stanza. Letteralmente. Un paio di tizi mi ha puntato
contro le armi, proprio dove sei ora. Di tutti i soldi che ho tirato su
quella notte, ne ho dovuti spendere il doppio per sistemare i danni.
Gente che beveva nelle camere, vomitava, spaccava tutto, un delirio. È
successo di peggio sull'Interstatale 10, più a nord. Il portiere
notturno al Days Inn, verso Ehrenberg, è stato pugnalato a morte. È
stato in quel momento che ho installato la recinzione di sicurezza, dopo
quello che era successo. Ora, appena inizia lo sfarfallio, spengo la
luce della scritta {CAMERE LIBERE} e rimango chiuso finché non è
finito.»

«E ascolti Duke» aggiunsi.

Sorrise. Entrammo per fare la registrazione. «Duke,» disse, «o Pops, o
Diz. Miles se sono dell'umore.» L'intimità dei veri fan nel chiamare i
morti per nome. «Niente di successivo al 1965, più o meno.» La reception
era illuminata da una luce tetra, tappezzata in modo approssimativo di
moquette e agghindata a tema vecchio West, ma dalla porta aperta sulla
stanza privata del proprietario - apparentemente dove viveva -
sgocciolava altra musica. Esaminò la carta di credito che gli diedi.

«Dottor Dupree,» lesse a voce alta porgendomi la mano, «sono Allen
Fulton. Sei diretto in Arizona?»

Gli raccontai che ero stato costretto a uscire dall'interstatale vicino
al confine.

«Non credo che ti andrà meglio sulla 10. Nelle notti come questa pare
che tutta Los Angeles voglia spostarsi a est. Come se lo sfarfallio
fosse una sorta di terremoto o maremoto.»

«Tra poco sarò di nuovo in strada.»

Mi porse una chiave. «Fatti prima un sonnellino. Ascolta il mio
consiglio.»

«Va bene la carta? Se preferisci contanti\ldots»

«Finché il mondo non finisce, la carta va bene quanto i contanti. E se
capiterà, dubito che avrò il tempo di pentirmene.»

Si mise a ridere. Io provai a sorridere.

Dieci minuti più tardi ero sdraiato, completamente vestito, su un letto
rigido in una stanza che puzzava di disinfettante al profumo di
pot-pourri e di aria condizionata eccessivamente umida, chiedendomi
perché non fossi rimasto in strada. Poggiai il telefono sul comodino,
chiusi gli occhi e mi addormentai all'istante.

ooo

E mi risvegliai dopo meno di un'ora, allarmato senza capirne la ragione.

Mi sedetti ed esaminai accuratamente la camera, tentando di riconoscere
le forme grigie nel buio secondo ciò che ricordavo. Alla fine la mia
attenzione si concentrò sul pallido rettangolo della finestra, sulla
tendina gialla che da quando ero entrato continuava a pulsare di luce.

Lo sfarfallio si era fermato.

E quel buio delicato avrebbe potuto rendere più facile dormire, ma
sapevo, essendoci già passato, che il sonno era fuori discussione. Lo
avevo imbrigliato per breve tempo, ma ormai se l'era data a gambe ed era
inutile far finta di niente.

Mi preparai un caffè con la piccola caffettiera di cortesia e ne bevvi
una tazza. Mezz'ora dopo, controllai di nuovo l'orologio. Due meno un
quarto. Piena notte. La zona dell'oggettività perduta. Avrei anche
potuto farmi una doccia e rimettermi in strada.

Mi vestii e percorsi il silenzioso vialetto di cemento fino all'ingresso
del motel prevedendo di fare scivolare la chiave nella fessura per la
posta; ma Fulton, il proprietario, era ancora sveglio, la luce del
televisore pulsava dalla sua stanza sul retro. Fece capolino quando
aprii la porta facendola cigolare.

Aveva un'aria strana. Un po' ubriaco o forse un po' strafatto. Batté le
palpebre finché mi riconobbe. «Dottor Dupree» borbottò.

«Mi dispiace disturbarti di nuovo. Devo rimettermi in viaggio. Grazie
dell'ospitalità, comunque.»

«Non mi devi dare spiegazioni» rispose. «Ti auguro il meglio e spero che
arriverai da qualche parte, prima dell'alba.»

«Me lo auguro anch'io.»

«Io lo sto guardando proprio adesso in TV.»

«Eh?»

All'improvviso non capivo di cosa stesse parlando.

«A volume basso. Non voglio svegliare Jody. Ti ho parlato di Jody? Mia
figlia. Ha dieci anni. Sua madre vive a La Jolla con un restauratore di
mobili. Jody passa l'estate con me. Qui in mezzo al deserto, che scherzo
del destino, vero?»

«Sì, be'\ldots»

«Ma non voglio svegliarla.» A un tratto sembrò incupirsi. «È sbagliato?
Lasciarla dormire tranquilla mentre succede? O finché può? O forse
dovrei svegliarla. Ora che ci penso, non le ha mai viste. Ha dieci anni.
Mai viste. Suppongo che sia la sua ultima opportunità.»

«Scusa, non sono sicuro di aver capito\ldots»

«Sono diverse, però. Non sono come me le ricordavo. Non che sia mai
stato una specie di esperto\ldots{} ma ai vecchi tempi, se passavi
abbastanza notti all'aperto, cominciavi a prenderci una certa
familiarità.»

«Familiarità con cosa?»

Batté gli occhi. «Le stelle» sussurrò.

ooo

Uscimmo a guardare il cielo vicino alla piscina vuota.

La piscina non veniva riempita da molto tempo. La polvere e la sabbia
avevano formato delle dune sul fondo e qualcuno si era firmato sulle
pareti con dei graffiti viola a caratteri balloon.

Il vento faceva sbatacchiare un'insegna d'acciaio ({NESSUN BAGNINO IN
	SERVIZIO}) contro le maglie della recinzione. Era un vento caldo e
proveniva da est.

Le stelle.

«Vedi?» disse. «Sono diverse. Non c'è nessuna delle vecchie
costellazioni. Tutto sembra come\ldots{} sparpagliato.»

Era stato l'effetto del passaggio di qualche miliardo di anni. Tutto
invecchiava, pure il cielo; l'intero universo tendeva alla massima
entropia, al disordine, alla casualità. Nel corso degli ultimi tre
miliardi di anni, la galassia in cui vivevamo era stata scossa in modo
massiccio da una violenza invisibile, che ne aveva fatto vorticare il
contenuto assieme a un ammasso galattico satellite (noto come M41 nei
vecchi cataloghi) fino a spargere le stelle per il cielo, in una distesa
senza senso. Era come guardare la mano rozza del tempo.

Fulton mi chiese: «Tutto a posto, dottor Dupree? Forse ti dovresti
sedere.»

Ero davvero troppo istupidito per stare in piedi. Mi sedetti sul cemento
rivestito di gomma facendo dondolare i piedi nella parte più bassa della
pendenza della piscina, continuando a fissare il cielo. Non avevo mai
visto niente di più bello né di più spaventoso.

«Solo poche ore al sorgere del Sole» mormorò Fulton con aria desolata.

Lì. Più a est, da qualche parte sull'Atlantico, il Sole doveva avere già
sfondato l'orizzonte. Volevo chiedergli cosa ne pensasse ma fui
interrotto da una vocina proveniente dall'oscurità vicino alla porta
della reception.

«Papà? Ti ho sentito parlare.»

Quella doveva essere Jody la figlia. Si avvicinò esitante. Indossava un
pigiama bianco e un paio di scarpe da ginnastica slacciate a proteggere
i piedi. Il suo viso era ampio, ordinario ma carino, e gli occhi
assonnati.

«Vieni qui, tesoro» disse Fulton. «Sali sulle mie spalle e dai
un'occhiata al cielo.»

Si arrampicò su, ancora confusa. Fulton rimase in piedi, le mani sulle
sue caviglie, sollevandola per avvicinarla il più possibile al buio
scintillante.

«Guarda,» esclamò sorridendo nonostante le lacrime che gli rigavano il
volto, «guarda là, Jody. Guarda tutto ciò che puoi, stanotte! Stanotte
potrai davvero vedere ogni cosa, dall'inizio alla fine di tutto.»

ooo

Tornai nella stanza per controllare se alla TV stessero trasmettendo
notiziari - Fulton mi aveva detto che la maggior parte dei telegiornali
delle stazioni via cavo erano ancora in onda.

Lo sfarfallio era finito un'ora prima. Era semplicemente svanito,
assieme alla membrana dello Spin. Lo Spin se n'era andato in silenzio
come era arrivato, niente fanfare, nessun rumore a parte uno
scricchiolio indistinto di elettricità statica, che sembrava provenire
dalla faccia soleggiata del pianeta.

Il Sole.

Tre miliardi e spiccioli di anni più vecchio di quando lo Spin lo aveva
sigillato. Provai a ricordare ciò che mi aveva detto Jase sulla
condizione del Sole a quel tempo. Funesta, senza dubbio; eravamo fuori
dalla zona vivibile; era risaputo. Sulla stampa l'immagine degli oceani
in ebollizione era stata argomento di discussione; ma eravamo già a quel
punto? Saremmo morti entro mezzogiorno oppure avremmo avuto fino alla
fine della settimana?

Aveva importanza?

Accesi il piccolo televisore nella stanza del motel e trovai una
trasmissione in diretta da New York City. Il panico, quello vero, non
era ancora dilagato. Troppe persone stavano ancora dormendo. Solo quelli
che avevano preceduto i pendolari mattutini e avevano visto le stelle
avevano potuto trarne le ovvie conclusioni. La troupe di questa
specifica redazione via cavo, in un sogno febbrile di eroismo
giornalistico, aveva sistemato sul tetto del Todt Hill a Staten Island
una telecamera puntata verso est. La luce era fioca, il cielo a oriente
si stava rischiarando ma era ancora sgombro. Un paio di conduttori
televisivi che riuscivano a malapena a reggere la tensione si leggevano
l'un l'altro alcuni bollettini appena ricevuti via fax.

Non c'era stato alcun collegamento intellegibile con l'Europa dalla fine
dello sfarfallio, dissero. Tutto ciò poteva essere dovuto alle
interferenze elettrostatiche, la luce solare non filtrata faceva
collassare i segnali provenienti dagli aerostati. Era troppo presto per
trarre conclusioni catastrofiche. «E come al solito,» disse uno dei
mezzibusti, «sebbene non ci sia alcun commento ufficiale, il consiglio è
quello di stare fermi dove siete e rimanere sintonizzati fino a quando
la situazione non sarà risolta. Non uscite se non per motivi
indispensabili.»

«Oggi più che mai,» il suo collega era d'accordo, «le persone vorranno
stare vicine alle proprie famiglie.»

Mi misi seduto sul bordo del materasso della camera del motel e
continuai a guardare il cielo ripreso dal notiziario fino al sorgere del
Sole.

La telecamera alta lo colse in un primo momento come uno strato di
nuvole cremisi che rasentavano l'orizzonte oleoso dell'Atlantico. Poi un
bordo di mezzaluna ribollente, e filtri che scorrevano sulle lenti per
bloccarne il bagliore.

Era difficile valutarne le dimensioni, ma il Sole sbucò fuori (non
proprio rosso ma di un arancio vermiglio, ammesso non fosse un effetto
artificioso della telecamera) e poi ancora più fuori, e continuò a
salire fino a librarsi su oceano, Queens, Manhattan, tanto grande da non
sembrare più un corpo celeste ma un enorme palloncino riempito di luce
ambrata.

Rimasi in attesa di nuovi ragguagli ma l'immagine rimase senza audio
finché non interruppero la ripresa e si collegarono con uno studio del
Midwest, il quartier generale di riserva della rete, da cui un altro
inviato, troppo poco tirato a lucido per essere un presentatore
abituale, proferì l'ennesima serie di avvertenze futili e infondate.
Spensi la TV.

Portai in macchina il mio equipaggiamento medico e la valigia.

Fulton e Jody uscirono dall'ufficio per salutarmi. D'improvviso erano
vecchi amici, spiacenti di vedermi andare via. Jody ora sembrava
spaventata. «Jody ha parlato con sua madre» mi spiegò Fulton. «Non credo
che sua madre avesse sentito delle stelle.»

Cercai di non immaginare quella sveglia di prima mattina, Jody che
telefonava dal deserto per dare alla madre una notizia che l'altra
avrebbe interpretato all'istante come la fine del mondo imminente. La
mamma di Jody che dà l'ultimo addio alla figlia e si sforza di non
spaventarla a morte, proteggendola dall'incombere della verità.

In quel momento, Jody era appoggiata alle costole di suo padre e Fulton
l'abbracciava, nient'altro che tenerezza tra loro.

«Devi proprio andare?» chiese Jody.

Risposi che dovevo.

«Perché se vuoi puoi restare. L'ha detto papà.»

«Il signor Dupree è un medico» le disse Fulton dolcemente.
«Probabilmente deve fare una visita a domicilio.»

«Esatto» confermai. «Devo proprio.»

ooo

Quella mattina, sulle corsie dell'autostrada dirette a est, successe
qualcosa di simile a un miracolo. Molte persone, in quelle che pensavano
fossero le loro ultime ore, diedero il peggio di sé. Era come se gli
sfarfallii fossero stati solo una prova generale per la fine del mondo.
Tutti quanti avevamo sentito le profezie: foreste in fiamme, caldo
rovente, mari trasformati in vapore vivo e ustionante. L'unico dubbio
era se ci sarebbe voluto un giorno, una settimana o un mese.

Perciò la gente prese a rompere vetrine e a rubare ciò che voleva,
qualunque cianfrusaglia la vita le avesse negato; gli uomini provavano a
violentare le donne e alcuni scoprivano che la perdita delle inibizioni
funzionava nei due sensi. Le vittime designate trovavano quindi una
forza inattesa in grado di cavare occhi e frantumare testicoli; vecchie
ruggini venivano risolte a colpi d'arma da fuoco e altri spari esplosi
solo per togliersi uno sfizio. I suicidi furono innumerevoli. Pensai a
Molly: se non fosse morta nel primo sfarfallio, a questo punto lo
sarebbe stata sicuramente. Addirittura, poteva essere morta compiaciuta
per il logico svolgimento del suo logico piano. E questo mi fece
piangere per lei per la prima volta in vita mia.

Ma vi furono anche isole di civiltà e atti di eroico altruismo.
L'Interstatale 10 al confine dell'Arizona ne fu un esempio.

Durante lo sfarfallio, vicino al ponte sul fiume Colorado c'era di
stanza un distaccamento della Guardia Nazionale. I soldati erano spariti
poco dopo la fine dello sfarfallio, forse richiamati, o solo disertori
che tornavano a casa. Senza di loro, il ponte sarebbe potuto diventare
un collo di bottiglia ingarbugliato e invalicabile.

Ma non fu così. Il traffico proseguì fluido in entrambi i sensi di
marcia. Una decina di civili, autoproclamatisi volontari e muniti di
robuste torce elettriche e bengala provenienti dai kit d'emergenza nei
bagagliai dei loro camion, si incaricarono di dirigere il traffico.
Perfino gli eterni ansiosi - quelli che volevano o dovevano viaggiare a
lungo prima dell'alba per raggiungere il Nuovo Messico, il Texas, forse
anche la Louisiana se il motore non si squagliava prima - sembrarono
rendersi conto di non avere alternative, nessun tentativo di passare
avanti avrebbe avuto successo e la pazienza era l'unica risorsa. Non so
quanto durò questo clima né quale confluenza di buone intenzioni e
casualità contribuì a crearlo. Forse si trattò di umana gentilezza o
forse solo di condizioni meteorologiche favorevoli: nonostante
l'imminente sciagura che ruggiva verso di noi da oriente, per assurdo la
notte fu \emph{gradevole}. Un cielo stupendo, limpido e disseminato di
stelle; una brezza frizzante, che si portò via il puzzo di gas di
scarico ed entrò dal finestrino dell'automobile delicata come il tocco
di una madre.

ooo

Pensavo di fare volontariato presso uno degli ospedali locali - al Palo
Verde di Blythe, che una volta avevo visitato per una conferenza, o
magari al La Paz Regional di Parker. Ma a quale scopo? Per ciò che stava
arrivando non c'era cura.

C'erano solo palliativi, morfina, eroina o la strada scelta da Molly,
sperando che gli armadietti dei medicinali non fossero già stati
saccheggiati.

E quello che aveva detto Fulton a Jody era fondamentalmente vero: dovevo
fare una visita a domicilio.

Era diventata la mia missione. E adesso sembrava proprio
donchisciottesca. Qualunque fosse il problema di Diane, non sarei
riuscito nemmeno a risolverlo. Allora perché andare da lei? Era una cosa
da fare alla fine del mondo, le mani impegnate non tremano, le menti
impegnate non vanno nel panico; ma tutto questo non spiegava l'urgenza,
il bisogno viscerale di vederla che mi aveva spinto in strada durante lo
sfarfallio e che, se possibile, in quel momento mi parve ancora più
forte.

Passata Blythe, passate le inquiete forche caudine dei negozi in
penombra e le zuffe che si formavano nelle stazioni di benzina
assediate, la strada si aprì. Il cielo era diventato di nuovo scuro e le
stelle brillavano. Mentre pensavo a tutto questo, squillò il telefono.

Rischiai quasi di finire fuori strada, frugando nella tasca e frenando,
e un'utilitaria che mi seguiva mi sorpassò stridendo.

«Tyler» disse Simon.

Prima che continuasse, lo interruppi: «Lasciami un numero da richiamare
prima di riagganciare o prima che s'interrompa la comunicazione. Così
posso raggiungerti.»

«Non posso farlo. Io\ldots»

«Stai chiamando dal tuo telefono personale o da quello di casa?»

«Una via di mezzo. È un cellulare, lo usiamo solo per le chiamate
locali. Ora ce l'ho io, ma Aaron lo porta con sé a volte, quindi\ldots»

«Non chiamerò, se proprio non devo.»

«Bene. Suppongo che non sia un problema.» Mi diede il numero. «Ma hai
visto il cielo, Tyler? Immagino di sì, dato che sei sveglio. È l'ultima
notte del mondo, o no?»

Pensai: ``Perché lo chiedi a me?'' Simon viveva negli \emph{ultimi
	giorni} da tre decenni, ormai. Avrebbe dovuto saperlo. «Dimmi di Diane»
ribattei.

«Voglio scusarmi per quella chiamata. Perché\ldots{} lo sai, per quello
che sta succedendo.»

«Come sta?»

«È quello che ti sto dicendo. Non importa.»

«È morta?»

Lunga pausa. Quando ricominciò a parlare sembrava offeso. «No. No, non è
morta. Non è quello il punto.»

«Sta forse fluttuando a mezz'aria in attesa dell'Estasi?»

«Non c'è bisogno di insultare la mia fede» mi rimproverò Simon stizzito.
(E io non potei fare a meno di interpretare la frase: ``la \emph{mia}
fede'', aveva detto, non la \emph{nostra} fede.)

«Perché, se non è così, forse ha bisogno di cure mediche. Sta ancora
male, Simon?»

«Sì. Ma\ldots»

«Quanto male? Quali sono i sintomi?»

«Manca solo un'ora al sorgere del Sole, Tyler. Sicuramente capisci cosa
significhi.»

«Non sono per nulla sicuro di ciò che significhi. E sono per strada,
posso arrivare al ranch prima dell'alba.»

«Ah\ldots{} no, non va bene\ldots{} no, io\ldots»

«Perché no? Se è la fine del mondo perché non dovrei essere lì?»

«Non capisci. Quello che sta succedendo non è solo la fine del mondo. È
la nascita di un mondo nuovo.»

«Quanto sta male esattamente? Posso parlarle?»

La voce di Simon si fece tremula. Un uomo al limite. Eravamo tutti al
limite. «Riesce appena a sussurrare. Non riesce a prendere fiato. È
debole. Ha perso molto peso.»

«Da quanto tempo sta così?»

«Non lo so. Voglio dire, è iniziato per gradi\ldots»

«Quando avete capito che era malata?»

«Qualche settimana fa. Anzi no, sarà passato qualche mese.»

«Ha ricevuto assistenza medica di qualche tipo?» Pausa. «Simon?»

«No.»

«Perché no?»

«Non sembrava necessario.»

«Non sembrava \emph{necessario}?»

«Il Pastore Dan non lo avrebbe permesso.»

Pensai: ``E tu hai detto al Pastore Dan di andarsene a fanculo?''

«Spero che abbia cambiato idea.»

«No\ldots»

«Perché altrimenti avrò bisogno del tuo aiuto per arrivare a lei.»

«Non farlo, Tyler. Non faresti del bene a nessuno.»

Stavo già cercando l'uscita, non ne ricordavo l'altezza, ma l'avevo
segnata sulla mappa. Era lontana dall'autostrada, verso una sorta di
palude secca come un osso, una strada deserta senza nome.

«Ha chiesto di me?» chiesi.

Silenzio.

«Simon? Ha chiesto di me?»

«Sì.»

«Dille che sarò lì al più presto possibile.»

«No, Tyler\ldots{} Tyler, ci sono dei problemi al ranch. Sul serio, non
puoi entrare qui.»

Problemi? «Pensavo stesse per nascere un nuovo mondo.»

«Nato nel sangue» rispose Simon.